\documentclass[11pt,a4paper]{article}

\usepackage{amsmath,amssymb,amsthm}
\usepackage{mathtools}
\usepackage{bm}
\usepackage{hyperref}
\usepackage[margin=1in]{geometry}
\usepackage{setspace}
\onehalfspacing

% Theorem environments
\newtheorem{theorem}{Theorem}[section]
\newtheorem{lemma}[theorem]{Lemma}
\newtheorem{proposition}[theorem]{Proposition}
\theoremstyle{definition}
\newtheorem{definition}[theorem]{Definition}
\newtheorem{example}[theorem]{Example}
\newtheorem{remark}[theorem]{Remark}

% Notation
\newcommand{\E}{\mathbb{E}}
\newcommand{\Var}{\mathrm{Var}}
\newcommand{\Cov}{\mathrm{Cov}}
\newcommand{\Corr}{\mathrm{Corr}}
\newcommand{\dpm}{m^*}
\newcommand{\dpv}{v^*}
\newcommand{\dpc}{c^*}
\newcommand{\dpr}{r^*}
\newcommand{\CE}{\mathrm{CE}}

\title{Covariance Decomposition for Quantity-Space Hyperpriors:\\
A First Pass}
\author{Technical Note}
\date{9 April 2026}

\begin{document}
\maketitle

\section{Motivation}

The published framework decomposes the variance of a single direct prior
$\dpv_k$ into a between-population component $\rho_k \dpv_k$ and a
within-population component $(1-\rho_k) \dpv_k$, then parameterizes the
heterogeneity confidence with $\kappa_k$. Each predictive quantity $k$
gets its own univariate decomposition, and the joint direct prior is
constructed by stacking independent marginals.

This is a deliberate simplification. In some applications, the researcher
has substantive beliefs about the joint distribution of two or more
predictive quantities. Two CEs computed for similar lotteries are an
obvious example: a researcher who believes the population is mostly
risk-averse should also believe the two CEs are positively correlated,
both within an individual (risk-averse subjects have low CEs for both
lotteries) and across the population (a more risk-averse population has
lower mean CEs for both). Stacking independent marginal direct priors
ignores this dependence.

This note develops a multivariate extension of the variance decomposition
that allows the researcher to specify a joint direct prior with nonzero
covariance between predictive quantities, and to decompose that joint
prior into a between-population and a within-population component using
the law of total covariance. I work through the simplest nontrivial case
($k = 2$ CEs that the researcher believes are positively correlated) and
sketch the open issues for the general case.

\section{Setup}

Let $v_i^1$ and $v_i^2$ be two predictive quantities for individual $i$
in an exchangeable population. The researcher specifies a joint direct
prior with means $\dpm_1, \dpm_2$, variances $\dpv_1, \dpv_2$, and
covariance $\dpc_{12}$ (or equivalently a correlation
$\dpr_{12} = \dpc_{12} / \sqrt{\dpv_1 \dpv_2}$). The hierarchical model
has individual-level distributions $p(v_i^k \mid \mu_k, V_k)$ for each
$k$, where $\mu_k$ is the population mean of quantity $k$ and $V_k$ is
the within-population variance, both treated as random under the
hyperprior.

The univariate framework matches two moments per quantity. The bivariate
extension matches the same two moments per quantity plus one cross-moment
($\dpc_{12}$), giving five constraints total. With a joint hyperprior on
$(\mu_1, \mu_2, V_1, V_2)$ that has enough free parameters, all five
constraints can be matched by moment-matching.

\section{The Law of Total Covariance}

The univariate decomposition rests on the law of total variance:
\begin{equation}
  \dpv_k
  = \E[\Var(v_i^k \mid \mu_k, V_k)] + \Var[\E(v_i^k \mid \mu_k, V_k)]
  = \E[V_k] + \Var[\mu_k].
\end{equation}
The first term is the expected within-population variance and the
second is the between-population variance.

The bivariate analogue is the law of total covariance:
\begin{equation}
  \label{eq:total_cov}
  \dpc_{12}
  = \E[\Cov(v_i^1, v_i^2 \mid \mu_1, \mu_2, V_1, V_2)]
  + \Cov[\E(v_i^1 \mid \cdot), \E(v_i^2 \mid \cdot)],
\end{equation}
which decomposes the total covariance between the two predictive
quantities into an expected within-population covariance and a
between-population covariance of population means:
\begin{equation}
  \label{eq:cov_decomp}
  \dpc_{12}
  = \underbrace{\E[\Cov_{\text{within}}(v_i^1, v_i^2 \mid \cdot)]}
        _{\text{within-pop covariance}}
  + \underbrace{\Cov[\mu_1, \mu_2]}_{\text{between-pop covariance}}.
\end{equation}

This decomposition is the natural multivariate generalization of
$\dpv_k = \E[V_k] + \Var[\mu_k]$. It says that two individuals from the
same population can have correlated quantities (the within term), and
that uncertainty about the population mean of one quantity can correlate
with uncertainty about the other (the between term).

\section{Two Substantively Distinct Sources of Correlation}

The two terms in~(\ref{eq:cov_decomp}) reflect distinct beliefs and the
researcher should be able to specify them separately.

\paragraph{Between-population covariance.}
This is the part of the joint uncertainty that comes from not knowing
which population the subjects are drawn from. If the researcher's belief
is that the population could be either ``mostly risk-averse'' (low mean
CE for both lotteries) or ``mostly risk-tolerant'' (high mean CE for
both), then $\Cov[\mu_1, \mu_2] > 0$. If the researcher believes the two
population means are independent (perhaps because they reflect different
preference channels), then $\Cov[\mu_1, \mu_2] = 0$.

\paragraph{Within-population covariance.}
This is the part of the joint uncertainty that comes from individual
heterogeneity even if the population means were known exactly. For two
CEs computed from the same individual under the same preference
parameters $(r, \phi, \eta)$, this within-individual covariance is
mechanically nonzero: a subject with high $r$ has low CEs for both
lotteries. The within-population covariance reflects the average
strength of this mechanical link across individuals.

\section{Parameterization Through Correlations}

Just as the univariate $\rho_k$ partitions $\dpv_k$ into between and
within components, the bivariate covariance can be parameterized through
two correlation-like quantities, one for each component
of~(\ref{eq:cov_decomp}).

Let $\rho^{\text{pop}}_{12}$ denote the correlation between the
population means $\mu_1$ and $\mu_2$, and let
$\rho^{\text{wn}}_{12}$ denote the expected within-population
correlation. Then
\begin{equation}
  \Cov[\mu_1, \mu_2]
  = \rho^{\text{pop}}_{12}\,\sqrt{\Var[\mu_1]\,\Var[\mu_2]}
  = \rho^{\text{pop}}_{12}\,\sqrt{\rho_1\,\dpv_1}\,\sqrt{\rho_2\,\dpv_2}
\end{equation}
and (under the simplifying assumption that the within-population
correlation does not depend on $V_1, V_2$)
\begin{equation}
  \E[\Cov_{\text{within}}(v_i^1, v_i^2 \mid \cdot)]
  \approx \rho^{\text{wn}}_{12}\,\sqrt{\E[V_1]\,\E[V_2]}
  = \rho^{\text{wn}}_{12}\,\sqrt{(1-\rho_1)\dpv_1}\,\sqrt{(1-\rho_2)\dpv_2}.
\end{equation}
Substituting into~(\ref{eq:cov_decomp}) and dividing through by
$\sqrt{\dpv_1 \dpv_2}$ gives the total correlation as a weighted
combination of the two component correlations:
\begin{equation}
  \label{eq:total_corr}
  \dpr_{12}
  = \rho^{\text{pop}}_{12}\sqrt{\rho_1 \rho_2}
  + \rho^{\text{wn}}_{12}\sqrt{(1-\rho_1)(1-\rho_2)}.
\end{equation}

The weights $\sqrt{\rho_1 \rho_2}$ and $\sqrt{(1-\rho_1)(1-\rho_2)}$
are the geometric means of the variance-partition ratios for the two
quantities. In the symmetric case $\rho_1 = \rho_2 = \rho$, these reduce
to $\rho$ and $1 - \rho$, recovering the univariate intuition.
Equation~(\ref{eq:total_corr}) is the bivariate generalization of the
total-variance decomposition: it tells the researcher how the elicited
total correlation $\dpr_{12}$ partitions into between- and
within-population components, given the chosen variance-partition ratios.

\section{Researcher Inputs in the Bivariate Case}

For two quantities, the researcher specifies:
\begin{itemize}
  \item Marginal direct prior moments $(\dpm_k, \dpv_k)$ for each
    quantity $k = 1, 2$.
  \item Variance partition ratios $\rho_1, \rho_2$ as in the univariate
    case.
  \item Heterogeneity confidence parameters $\kappa_1, \kappa_2$ as in
    the univariate case.
  \item One additional input governing the joint structure: either the
    between-population correlation $\rho^{\text{pop}}_{12}$, or the
    within-population correlation $\rho^{\text{wn}}_{12}$, or the total
    correlation $\dpr_{12}$ (with a choice of how to partition it).
\end{itemize}
This is one extra researcher input per pair of quantities. With $p$
quantities the joint specification adds $\binom{p}{2}$ correlation
inputs in addition to the $4p$ univariate inputs already required.

A natural way to organize these is to elicit the total correlation
$\dpr_{12}$ (the easiest object for a researcher to reason about) and
then choose how to split it between the two components. Two practical
defaults are:
\begin{itemize}
  \item \textit{Match the variance partition.} Set
    $\rho^{\text{pop}}_{12} = \rho^{\text{wn}}_{12} = \dpr_{12}$.
    This is consistent with the univariate ratio:
    the same fraction of the joint uncertainty is allocated to the
    between-population component as in the marginal case.
  \item \textit{Concentrate within.} Set
    $\rho^{\text{pop}}_{12} = 0$ and let the within-population
    correlation absorb all of $\dpr_{12}$. This reflects a belief that
    the two quantities co-vary because they are functions of the same
    individual parameters, but not because the population mean is
    uncertain in a way that links them.
\end{itemize}
The first default is the most parsimonious starting point. The second
is appropriate when the researcher believes $\Cov[\mu_1, \mu_2] = 0$ on
substantive grounds.

\section{A Concrete Example: Two Positively Correlated CEs}

Consider two binary lotteries from the running example: $L_1$ paying
$(\$10, \$50)$ with probability $0.25$ on the high outcome, and $L_2$
paying $(\$20, \$40)$ with probability $0.5$. Both lotteries have
expected value $\$20$. Suppose the researcher elicits the following
direct prior beliefs:
\begin{align*}
  \dpm_1 &= 18.5,
  & \dpv_1 &= 2.08,
  & \dpr_{12} &= 0.6, \\
  \dpm_2 &= 28.5,
  & \dpv_2 &= 2.05.
\end{align*}
The marginal moments come from the published application. The joint
input is the total correlation $\dpr_{12} = 0.6$, reflecting a belief
that subjects with low $\CE_1$ also tend to have low $\CE_2$, both
within the population (because the mechanical channel through $r$ is
strong) and between possible populations (because the population is
either mostly risk-averse or mostly risk-tolerant).

Take the symmetric variance-partition $\rho_1 = \rho_2 = 0.2$ used in
the published application. Under the parsimonious default
$\rho^{\text{pop}}_{12} = \rho^{\text{wn}}_{12} = 0.6$,
equation~(\ref{eq:total_corr}) gives
\[
  \dpr_{12} = 0.6 \cdot 0.2 + 0.6 \cdot 0.8 = 0.6,
\]
as required. The between-population covariance is
\[
  \Cov[\mu_1, \mu_2]
  = 0.6 \cdot \sqrt{0.2 \cdot 2.08 \cdot 0.2 \cdot 2.05}
  \approx 0.246,
\]
and the expected within-population covariance is
\[
  \E[\Cov_{\text{within}}]
  = 0.6 \cdot \sqrt{0.8 \cdot 2.08 \cdot 0.8 \cdot 2.05}
  \approx 0.984.
\]
These two pieces sum to $\dpc_{12} = 0.6 \cdot \sqrt{2.08 \cdot 2.05}
\approx 1.230$, as required by~(\ref{eq:cov_decomp}).

\section{Implementing the Joint Hyperprior}

The decomposition above tells the researcher how the elicited
correlation partitions into between- and within-population components.
The remaining problem is to construct a joint hyperprior over
$(\mu_1, \mu_2, V_1, V_2)$ whose induced moments match these targets.
For the bivariate case, this requires choosing:
\begin{itemize}
  \item A joint distribution for $(\mu_1, \mu_2)$ with specified means,
    variances, and correlation $\rho^{\text{pop}}_{12}$.
  \item A joint distribution for the individual-level draws
    $(v_i^1, v_i^2)$ given $(\mu_1, \mu_2, V_1, V_2)$ that delivers the
    target within-population correlation.
\end{itemize}

For Normal-Inverse-Gamma hierarchies, both pieces are straightforward:
$(\mu_1, \mu_2)$ can be jointly Normal with the required correlation,
and the individual-level draws can be jointly Normal with a correlation
matrix encoding $\rho^{\text{wn}}_{12}$. The variances $V_1, V_2$ can
be marginally Inverse-Gamma and either independent or coupled through a
copula.

For the Beta-Beta hierarchy used in the application, the joint
construction is more involved. Bivariate Beta distributions on the
unit square exist (Sarmanov, FGM, copula constructions) but the
moment-matching from $(\rho^{\text{pop}}_{12}, \rho^{\text{wn}}_{12})$
to the parameters of the joint distribution is no longer closed-form.
A practical approach is to sample-based moment-matching: simulate from
the candidate joint hyperprior, compute the empirical correlations, and
adjust the hyperparameters until they match the targets. This
sacrifices the closed-form solution of the univariate framework but
preserves the interpretability of the inputs.

\section{Open Issues}

Three issues remain for the general case.

\textit{Higher-order joint moments.} The decomposition above matches
first and second moments only. The bivariate generalization of $\kappa$
would require parameterizing the joint variance-of-variances in some
interpretable way, which is not obvious.

\textit{Feasibility regions.} For the univariate Beta-Beta hierarchy
the feasibility region is already a proper subset of $(0,1) \times
(0,\infty)$. With joint inputs the feasibility region becomes
higher-dimensional and may shrink considerably. Mapping out feasible
$(\rho^{\text{pop}}_{12}, \rho^{\text{wn}}_{12})$ pairs for a given
marginal specification is an open problem.

\textit{Beyond pairs.} For $p > 2$ the joint specification involves a
$p \times p$ correlation matrix at each level, with positive
semidefiniteness constraints linking the entries. Eliciting and
verifying these constraints is the main practical barrier to scaling
beyond bivariate joint priors.

\section{Summary}

The univariate variance decomposition extends naturally to a bivariate
covariance decomposition via the law of total covariance. The
researcher specifies one extra input per pair of quantities (a total
correlation $\dpr_{12}$, or equivalently a between-population
correlation), and the decomposition partitions it into between- and
within-population components using the same variance-partition ratios
$\rho_1, \rho_2$ already required by the univariate framework.

The implementation problem is the main remaining gap. For
Normal-Inverse-Gamma hierarchies it is straightforward; for Beta-Beta
and other non-Gaussian families it requires either sample-based
moment-matching or new analytical work on bivariate distributions with
prescribed marginals and correlation. The extension to $p > 2$ is
conceptually clean but practically demanding, mainly because of the
positive-semidefiniteness constraints on the joint correlation matrix.

\end{document}
